\begin{abstract}
\noindent
A deterioração clínica de pacientes internados em enfermarias é frequentemente detectada de forma tardia, em razão do monitoramento intermitente dos sinais vitais. Este trabalho propõe um sistema embarcado de beira de leito que monitora continuamente os sinais vitais, calcula o escore de alerta precoce NEWS2 na borda e prediz sua evolução por meio de aprendizado de máquina. O dispositivo é construído sobre uma Raspberry Pi 3B+ com sistema operacional Linux customizado gerado pelo Buildroot, integrando sensores de fotopletismografia e de temperatura, interface gráfica em Qt e periféricos de alerta. Para a predição, são comparados três modelos de vieses indutivos distintos (regressão logística, XGBoost e LSTM), treinados sobre uma base pública de pacientes hospitalizados. Como resultados preliminares, o \textit{hardware} foi definido e adquirido, os pacotes do sistema operacional embarcado foram definidos no Buildroot e a base de dados que fundamenta o modelo preditivo foi normalizada e preparada para o treinamento. Os avanços confirmam a viabilidade da solução, cuja finalização e validação serão conduzidas na etapa seguinte do trabalho.
\\\textbf{\keyword{Palavras-chave: }} NEWS2. Deterioração clínica. Sistemas embarcados. Aprendizado de máquina. Computação em borda.
\end{abstract}

\newpage
\selectlanguage{english}
\begin{abstract}
\noindent
The clinical deterioration of patients in hospital wards is often detected late, owing to the intermittent monitoring of vital signs. This work proposes a bedside embedded system that continuously monitors vital signs, computes the NEWS2 early warning score at the edge, and predicts its evolution through machine learning. The device is built on a Raspberry Pi 3B+ running a custom embedded Linux generated with Buildroot, integrating photoplethysmography and temperature sensors, a Qt graphical interface, and alert peripherals. For the prediction, three models with distinct inductive biases (logistic regression, XGBoost, and LSTM) are compared, trained on a public dataset of hospitalized patients. As preliminary results, the hardware was defined and acquired, the embedded operating system packages were defined in Buildroot, and the dataset underlying the predictive model was normalized and prepared for training. These advances confirm the feasibility of the solution, whose completion and validation will be carried out in the next stage of the work.
\\\textbf{\keyword{Keywords: }} NEWS2. Clinical deterioration. Embedded systems. Machine learning. Edge computing.
\end{abstract}
